\documentclass[10pt,a4paper]{amsart}
\usepackage[margin=19mm]{geometry}
\usepackage{amsmath,amssymb,mathtools}
\usepackage[scr=rsfs,cal=euler]{mathalfa}
\usepackage{xfrac}
\usepackage[unicode,hidelinks,bookmarks=false]{hyperref}
\newcommand{\ie}{i.e.\ }
\newcommand{\cf}{cf.\ }
\newcommand{\resp}{respectively\ }
\newcommand{\Subsection}{\subsection}
\providecommand{\nobreakdash}{\nobreak\hbox{-}}
\setlength{\emergencystretch}{2em}
\multlinegap0pt
\usepackage{bm}
\usepackage{bbm}
\usepackage{dsfont}
\usepackage{xspace}
\usepackage{cancel}
\usepackage{mathtools}
\usepackage{amsthm}
\makeatletter
\@ifundefined{thm}{\newtheorem{thm}{Thm}}{}
\@ifundefined{assumption}{\newtheorem{assumption}{Assumption}}{}
\@ifundefined{prop}{\newtheorem{prop}[thm]{Proposition}}{}
\makeatother
\newcommand{\B}{\mathcal{B}}
\newcommand{\R}{\mathbb{R}}
\title[Slowly oscillating periodic solutions in a nonlinear Volterr]{Slowly oscillating periodic solutions in a nonlinear Volterra equation with non-symmetric feedback}
\author{Quentin Griette, Franco Herrera}
\date{Source: arXiv:2506.09564v1 (CC BY 4.0)}
\begin{document}
\maketitle

\noindent\emph{Source and licence.} Slowly oscillating periodic solutions in a nonlinear Volterra equation with non-symmetric feedback. Version arXiv:2506.09564v1 (CC BY 4.0), distributed under the Creative Commons Attribution 4.0 International licence, as recorded in the arXiv licence metadata for this exact version and shown on the abstract page https://arxiv.org/abs/2506.09564v1. Full rights notice: licenses/arxiv-2506.09564-CC-BY-4.0.txt.\par\smallskip

\noindent\textit{Editorial scope.} Counted unit: the Prop whose source label is \texttt{prop:estimates1} in the article (source0.tex lines 672--681, source label \texttt{prop:estimates1}), delivered together with the article's own complete original proof of it (source0.tex lines 682--755, one \texttt{proof} environment of 74 lines). No mathematics is written, repaired or re-derived: statement and proof are the article's own text line for line. Required context retained uncounted: as:f (lines 168--178). Every definition, display and label of those lines is retained unchanged. Omitted from this package and not redistributed: 1--167, 179--671, 756--1456. Nothing inside a retained excerpt is omitted or altered: every retained body is delivered exactly as the article's lines are written, so no omission receipt of any kind accompanies this package. Citations: the retained excerpts carry no citation command at all, so no bibliography entry of the article is transcribed and no reference is redistributed. Reference adaptation: none was needed and none was made; every reference inside the delivered unit resolves inside it, so no unresolved reference or citation is produced. Printed slips kept literal: no printed word, symbol, hypothesis or number of the counted statement or of its proof was altered. Layout and class: the article's own class is \texttt{article} with packages including graphicx, inputenc, amsmath, amssymb, amsthm, mathrsfs, xcolor, setspace, geometry, enumitem, tikz, pstricks, pst-all, verbatim; the wrapper is the lane's pinned \texttt{amsart} editorial wrapper at 10pt on A4, which restates the theorem-like environments the retained text needs and adds the article's own macro definitions that the retained text needs (\texttt{\textbackslash B}, \texttt{\textbackslash R}), with extra wrapper packages (bm, bbm, dsfont, xspace, cancel, mathtools). The delivered numbering is the unit's own. Assets: no retained span contains a figure environment, a graphics-inclusion command, a PDF-page inclusion, a table or a TikZ picture; no image file is copied into this package, the package declares no asset dependency and no image file is redistributed with it. Licence: version arXiv:2506.09564v1 of this article is distributed under the Creative Commons Attribution 4.0 International licence, as recorded in the arXiv licence metadata for this exact version and shown on the abstract page https://arxiv.org/abs/2506.09564v1. A full rights notice is delivered at licenses/arxiv-2506.09564-CC-BY-4.0.txt. Characters: every retained excerpt is pure ASCII, so the delivered engine, XeLaTeX with its default Latin Modern text font, reports no missing character.
\begin{assumption}\label{as:f}
    $f\colon \R\to\R$ is a continuous function with $f'(0)<-1$, and satisfies the negative-feedback condition $xf(x)<0$. We let $\kappa_1$ and $\kappa_2$ be respectively the first negative and first positive solution of the equation $f(x)=-x$:
    \begin{align*}
	\kappa_1&:=\inf\{x<0\,:\, f(z)>-z, \quad  \,^\forall z\in [x, 0)\}, \\ 
	\kappa_2&:=\sup\{x>0\,:\, f(z)<-z,\quad  \,^\forall z\in (0, x]\}. 
    \end{align*}
    Furthermore, the following conditions hold
\[
    \limsup_{x\to +\infty} f(x)<0<\liminf_{x\to-\infty} f(x), \quad \text{and} \quad \liminf_{x\to \pm \infty} \frac{f(x)}{x} > -1.
\]
\end{assumption}

\begin{prop}[Estimates for $\tau\in(0,\varepsilon/2{]}$]\label{prop:estimates1}
Let $b\in\B^\alpha$ with $\tau\in(0,\varepsilon/2]$. Then, the extension $x_b$ verifies the following estimates:
\begin{align}
    x_b(t) &\leq \lambda_0 \left( 1+\frac{\delta_0}{f'(0)} \right) \min_{\hat{\tau}\in [\tau,\varepsilon/2]} \alpha \varphi_0^{\hat{\tau}} (t), \quad t\in[0,1-\tau-\varepsilon/2]; \label{estimate1} \\ 
    \frac{x_b(t)}{\alpha \varphi_0^{1-z_1(b)}(t)} &>1, \quad t\in(1-\tau-\varepsilon/2,1-\tau+\varepsilon/2) \setminus \{ z_1(b) \}; \label{stepness} \\     
    x_b(t) &\geq \alpha\lambda_0 \left( 1+\frac{\delta_0}{f'(0)} \right) \varphi_0^\tau (t), \quad t\in[1-\tau+\varepsilon/2, 2-\tau-\varepsilon];\label{estimate2} \\
    x_b(t) &\geq \alpha\lambda_0 \left( 1+\frac{\delta_0}{f'(0)} \right)\varphi_0^{1-z_1(b)}(t), \quad t\in(2-\tau-\varepsilon,z_1(b)+1-\varepsilon/2]; \label{estimate3} \\ 
    \frac{x_b(t)}{\alpha \varphi_0^{-z_2(b)}(t)} &>1, \quad t\in(z_1(b)+1-\varepsilon/2,z_1(b)+1+\varepsilon/2) \setminus\{ z_2(b) \} \label{stepness2}.
\end{align}
\end{prop}

\begin{proof}
To prove \eqref{estimate1} let $t\in [0,1-\tau-\varepsilon/2]$. Using Assumption \ref{as:f} and the previous remark we obtain
\begin{align*}
    x_b(t) &= \frac{1}{\varepsilon}\int_{t-1-\varepsilon/2}^{t-1+\varepsilon/2} f(b(s)) \\
    &\leq \frac{1}{\varepsilon}\int_{t-1-\varepsilon/2}^{t-1+\varepsilon/2} f\left( \max_{\hat{\tau}\in[\tau,\varepsilon/2]} \alpha \varphi_0^{\hat{\tau}} (s) \right) ds \\
    &\leq \frac{f'(0)+\delta_0}{\varepsilon}\int_{t-1-\varepsilon/2}^{t-1+\varepsilon/2} \max_{\hat{\tau}\in[\tau,\varepsilon/2]} \alpha \varphi_0^{\hat{\tau}} (s) ~ ds \\
    &\leq \alpha (f'(0) + \delta_0)\frac{\lambda_0}{f'(0)} \min_{\hat{\tau}\in[\tau,\varepsilon/2]} \varphi_0^{\hat{\tau}}(t),
\end{align*}
and thus \eqref{estimate1} follows. In addition, this inequality implies that
\begin{equation}\label{eq:inter-estimate}
    x_b(t) \leq \alpha\varphi_0^\tau(t), \quad \forall t\in[0,1-\tau-\varepsilon/2],
\end{equation}
due to the condition on $\delta_0$.

To prove \eqref{stepness} let $t\in(1-\tau-\varepsilon/2,1-\tau+\varepsilon/2)$. From the previous inequalities we have
\begin{align*}
    x_b'(t) &= \frac{1}{\varepsilon} [ f(x_b(t-1+\varepsilon/2)) - f(x_b(t-1-\varepsilon/2)) ] \\
    &\geq \frac{1}{\varepsilon} [f(\alpha\varphi_0^\tau(t-1+\varepsilon/2)) - f(\alpha\varphi_0^\tau(t-1-\varepsilon/2))] \\
    &\geq \alpha\frac{f'(0)+\delta_0}{\varepsilon} [\varphi_0^\tau(t-1+\varepsilon/2) - \varphi_0^\tau(t-1-\varepsilon/2)] \\
    &= \alpha (f'(0)+\delta_0)\frac{\lambda_0}{f'(0)} (\varphi_0^\tau)'(t) \\
    &\geq \alpha\pi\lambda_0\left( 1+\frac{\delta_0}{f'(0)} \right) \cos\left( \frac{\pi\varepsilon}{2} \right) \\
    &>\alpha\pi.
\end{align*}
This estimate on the derivative implies that $x_b'(t) < \alpha(\varphi_0^\eta)'(t)$ for all $\eta\in \R$ and $t$ in this interval, thus the functions $x_b$ and $\alpha\varphi_0^\eta$ intersects at most once, and therefore we deduce \eqref{stepness}.

For \eqref{estimate2}, let $t\in[1-\tau+\varepsilon/2,2-\tau-\varepsilon]$. As in the proof of \eqref{estimate1}, and using \eqref{eq:inter-estimate}, we have
\begin{align*}
    x_b(t) &= \frac{1}{\varepsilon} \int_{t-1-\varepsilon/2}^{t-1+\varepsilon/2} f(x_b(s)) ds \\
    & \geq \frac{1}{\varepsilon} \int_{t-1-\varepsilon/2}^{t-1+\varepsilon/2} f(\alpha\varphi_0^\tau(s)) ds \\
    &\geq \frac{f'(0)+\delta_0}{\varepsilon} \int_{t-1-\varepsilon/2}^{t-1+\varepsilon/2} \alpha\varphi_0^\tau(s) ds \\
    &=\alpha\lambda_0\left( 1+\frac{\delta_0}{f'(0)} \right)\varphi_0^\tau(t).
\end{align*}

Subsequently, for \eqref{estimate3} consider $t\in(2-\tau-\varepsilon,z_1(b)+1-\varepsilon/2]$, and observe that
\begin{align*}
    x_b(t) &= \frac{1}{\varepsilon} \left( \int_{t-1-\varepsilon/2}^{1-\tau-\varepsilon/2} f(x_b(s)) ds + \int_{1-\tau-\varepsilon/2}^{t-1+\varepsilon/2} f(x_b(s)) ds \right) \\
    &\geq \frac{1}{\varepsilon} \left( \int_{t-1-\varepsilon/2}^{1-\tau-\varepsilon/2} f\left( \alpha\lambda_0\left( 1+\frac{\delta_0}{f'(0)}\right)\varphi_0^\tau(s) \right) ds + \int_{1-\tau-\varepsilon/2}^{t-1+\varepsilon/2} f(\alpha\varphi_0^{1-z_1(b)}(s)) ds \right) \\
    &\geq \alpha\frac{f'(0)+\delta_0}{\varepsilon} \left( \int_{t-1-\varepsilon/2}^{1-\tau-\varepsilon/2} \lambda_0\left( 1+\frac{\delta_0}{f'(0)}\right)\varphi_0^\tau(s) ds + \int_{1-\tau-\varepsilon/2}^{t-1+\varepsilon/2} \varphi_0^{1-z_1(b)}(s) ds \right).
\end{align*}

Now, if $z_1(b)\leq 1-\tau$, it is easy to see that $\varphi_0^{\tau} \leq \varphi_0^{1-z_1(b)}$ in the first interval of integration. On the other hand, if $z_1(b)>1-\tau$ then $\lambda_0(1+\delta_0/f'(0))\varphi_0^\tau \leq \varphi_0^{1-z_1(b)}$ in the aforementioned interval. In fact, due to the shape of the functions involved, it is straightforward to check they intersect at most once in $[1/2-\tau, 1-\tau]$. Then, since
\[
    \lambda_0\left( 1+\frac{\delta_0}{f'(0)} \right) \varphi_0^\tau(1-\tau) = 0 >\varphi_0^{1-z_1(b)}(1-\tau),
\]
to get the claim it is enough to ensure that
\[
    \lambda_0\left( 1+\frac{\delta_0}{f'(0)} \right)\varphi_0^{\tau}(1-\tau-\varepsilon/2)\leq \varphi_0^{1-z_1(b)}(1-\tau-\varepsilon/2).
\]
This last inequality is equivalent to
\[
    \lambda_0\left( 1+\frac{\delta_0}{f'(0)} \right) \geq \frac{\varphi_0(1-z_1(b)-\tau-\varepsilon/2)}{\varphi_0(-\varepsilon/2)},
\]
and since $z_1(b)\in(1-\tau,1-\tau+\varepsilon/2)$, the right-hand side is bounded above by $\frac{\varphi_0(\varepsilon)}{\varphi_0(\varepsilon/2)}< 2$ and this last inequality holds. Consequently
\begin{align*}
    x_b(t) &\geq \alpha\frac{f'(0)+\delta_0}{\varepsilon} \int_{t-1-\varepsilon/2}^{t-1+\varepsilon/2} \varphi_0^{1-z_1(b)}(s) ds \\
    & = \alpha\lambda_0\left( 1+\frac{\delta_0}{f'(0)} \right) \varphi_0^{1-z_1(b)}(t),
\end{align*}
and thus \eqref{estimate3} follows.

Finally, for \eqref{stepness2} we use an analogous steepness argument as for \eqref{stepness}. Let $t\in (z_1(b)+1-\varepsilon/2, z_1(b)+1+\varepsilon/2)$, then we know that
\[
    x_b'(t) = \frac{1}{\varepsilon} [ f(x_b(t-1+\varepsilon/2)) - f(x_b(t-1-\varepsilon/2)) ].
\]
We divide the proof in two cases to use the different information about $x_b$, especially as the estimates change at $1-\tau-\varepsilon/2$. If $t-1-\varepsilon/2>1-\tau-\varepsilon/2$ then \eqref{stepness} and \eqref{estimate2} yield
\begin{align*}
    x_b'(t)& \leq \frac{1}{\varepsilon} \left[ f\left( \alpha\lambda_0\left( 1+\frac{\delta_0}{f'(0)} \right) \varphi_0^\tau(t-1+\varepsilon/2) \right) - f(\alpha\varphi_0^{1-z_1(b)}(t-1-\varepsilon/2)) \right] \\
    & \leq \alpha\frac{f'(0)+\delta_0}{\varepsilon} \left[ \lambda_0\left( 1+\frac{\delta_0}{f'(0)} \right) \varphi_0^\tau(t-1+\varepsilon/2) - \varphi_0^{1-z_1(b)}(t-1-\varepsilon/2) \right],
\end{align*}
and in the same fashion as for the previous statement, we deduce that $\lambda_0(1+\delta_0/f'(0))\varphi_0^\tau \geq \varphi_0^{1-z_1(b)}$ at $t-1+\varepsilon/2$, so that
\begin{align*}
    x_b'(t) &\leq \alpha\frac{f'(0)}{\varepsilon}[\varphi_0^{1-z_1(b)}(t-1+\varepsilon/2) - \varphi_0^{1-z_1(b)}(t-1-\varepsilon/2)] <-\alpha\pi,
\end{align*}
and we conclude just as for \eqref{stepness}. The complementary case when $t-1-\varepsilon/2\leq 1-\tau-\varepsilon/2$ is analogous.
\end{proof}

\end{document}
